\section{The finiteness theorem for proper morphisms}
\label{section:III.3}

\subsection{The dévissage lemma}
\label{subsection:III.3.1}

\begin{definition}[3.1.1]
\label{III.3.1.1}
Let $K$ be an abelian category. We say that a subset $K'$ of
the set of objects of $K$ is \emph{exact} if $0\in K'$ and if, for every
exact sequence
\[
0\longrightarrow A'\longrightarrow A\longrightarrow A''\longrightarrow0
\]
in $K$ such that two of the objects $A$, $A'$, $A''$ are in $K'$, the
third is also in $K'$.
\end{definition}

\begin{theorem}[3.1.2]
\label{III.3.1.2}
Let $X$ be a Noetherian prescheme; we denote by $K$ the abelian category
of coherent $\mathcal O_X$-modules. Let $K'$ be an exact subset of
$K$, and $X'$ a closed subset of the underlying space of $X$. Suppose that
for every irreducible closed subset $Y$ of $X'$, with generic point $y$, there
exists an $\mathcal O_X$-module $\mathcal G\in K'$ with support $Y$ such that $\mathcal G_y$ is
a $k(y)$-vector space of dimension $1$. Then every coherent $\mathcal O_X$-module
with support contained in $X'$ belongs to $K'$ (and in particular,
if $X'=X$, we have $K'=K$).
\end{theorem}

\begin{proof}
Consider the following property $P(Y)$ of a closed subset $Y$ of $X'$:
every coherent $\mathcal O_X$-module with support contained in $Y$ belongs to
$K'$. By virtue of the principle of Noetherian induction (0\textsubscript{I},
2.2.2), we see that we are reduced to proving that \emph{if $Y$ is a closed
subset of $X'$ such that the property $P(Y')$ is true for every closed
subset $Y'$ of $Y$ distinct from $Y$, then $P(Y)$ is true}.

So let $\mathcal F\in K$ have support contained in $Y$, and let us prove that
$\mathcal F\in K'$. We again denote by $Y$ the reduced closed sub-prescheme of
$X$ having $Y$ as underlying space (I, 5.2.1); it is defined by a coherent
ideal $\mathcal J$ of $\mathcal O_X$. We know (I, 9.3.4) that there exists an
integer $n>0$ such that $\mathcal J^n\mathcal F=0$; for $1\leq k\leq n$ we
therefore have an exact sequence
\[
0\longrightarrow
\mathcal J^{k-1}\mathcal F/\mathcal J^k\mathcal F
\longrightarrow \mathcal F/\mathcal J^k\mathcal F
\longrightarrow \mathcal F/\mathcal J^{k-1}\mathcal F
\longrightarrow0
\]
of coherent $\mathcal O_X$-modules (0\textsubscript{I}, 5.3.6 and 5.3.3);
since $K'$ is exact, we see, by induction on $k$, that it suffices to show
that each of the
$\mathcal F_k=\mathcal J^{k-1}\mathcal F/\mathcal J^k\mathcal F$ is in
$K'$. We are therefore reduced to proving that $\mathcal F\in K'$ under the
additional hypothesis that $\mathcal J\mathcal F=0$; this amounts to saying that
$\mathcal F=j_*(j^*(\mathcal F))$, where $j$ is the canonical injection
$Y\to X$. We now distinguish two cases:
\begin{enumerate}
\item[a)] $Y$ is \emph{reducible}. Let $Y=Y'\cup Y''$, where $Y'$ and $Y''$
are closed subsets of $Y$ distinct from $Y$; let $Y'$,
$Y''$ also denote the reduced closed sub-preschemes of $X$ having as underlying
spaces $Y'$, $Y''$ respectively, which are defined respectively by
the coherent ideals $\mathcal J'$, $\mathcal J''$ of $\mathcal O_X$. Set
\[
\mathcal F'=\mathcal F\otimes_{\mathcal O_X}
(\mathcal O_X/\mathcal J'),\qquad
\mathcal F''=\mathcal F\otimes_{\mathcal O_X}
(\mathcal O_X/\mathcal J'').
\]
The canonical homomorphisms $\mathcal F\to\mathcal F'$,
$\mathcal F\to\mathcal F''$ therefore define a homomorphism
$u:\mathcal F\to\mathcal F'\oplus\mathcal F''$. Let us show that for every
$z\notin Y'\cap Y''$, the homomorphism
$u_z:\mathcal F_z\to\mathcal F'_z\oplus\mathcal F''_z$ is \emph{bijective}.
Indeed, we have $\mathcal J'\cap\mathcal J''=\mathcal J$, since the question is
local and

\oldpage[III]{116}
the preceding equality follows from (I, 5.2.1 and 1.1.5); if $z\notin Y''$, we
therefore have $\mathcal J'_z=\mathcal J_z$, whence $\mathcal F'_z=\mathcal F_z$ and
$\mathcal F''_z=0$, which establishes our assertion in this case; we argue
in the same way if $z\notin Y'$. Consequently, the kernel and the cokernel of $u$,
which are in $K$ (0\textsubscript{I}, 5.3.4), have their supports in $Y'\cap Y''$,
and are therefore in $K'$ by hypothesis; for the same reason, $\mathcal F'$ and
$\mathcal F''$ are in $K'$, hence so is $\mathcal F'\oplus\mathcal F''$,
since $K'$ is exact. The conclusion then follows from consideration of the
two exact sequences
\[
0\longrightarrow\operatorname{Im}u\longrightarrow
\mathcal F'\oplus\mathcal F''\longrightarrow
\operatorname{Coker}u\longrightarrow0,
\]
\[
0\longrightarrow\operatorname{Ker}u\longrightarrow\mathcal F
\longrightarrow\operatorname{Im}u\longrightarrow0
\]
and from the hypothesis that $K'$ is exact.

\item[b)] $Y$ is irreducible, and consequently the sub-prescheme $Y$ of $X$ is
\emph{integral}. If $y$ is its generic point, we have
$(\mathcal O_Y)_y=k(y)$, and since $j^*(\mathcal F)$ is a coherent
$\mathcal O_Y$-module, $\mathcal F_y=(j^*(\mathcal F))_y$ is a
$k(y)$-vector space of finite dimension $m$. By hypothesis, there is a
coherent $\mathcal O_X$-module $\mathcal G$ (necessarily with support $Y$)
such that $\mathcal G_y$ is a $k(y)$-vector space of dimension $1$.
Consequently, there is a $k(y)$-isomorphism
$(\mathcal G_y)^m\xrightarrow{\sim}\mathcal F_y$, which is also an
$\mathcal O_Y$-isomorphism, and since $\mathcal G^m$ and $\mathcal F$ are
coherent, there exist an open neighbourhood $W$ of $y$ in $X$ and an
isomorphism $\mathcal G^m|W\xrightarrow{\sim}\mathcal F|W$
(0\textsubscript{I}, 5.2.7). Let $\mathcal H$ be the graph of this isomorphism,
which is a coherent sub-$(\mathcal O_X|W)$-module of
$(\mathcal G^m\oplus\mathcal F)|W$, canonically isomorphic to
$\mathcal G^m|W$ and to $\mathcal F|W$; there therefore exists a coherent
sub-$\mathcal O_X$-module $\mathcal H_0$ of $\mathcal G^m\oplus\mathcal F$
inducing $\mathcal H$ on $W$ and $0$ on $X-Y$, since $\mathcal G^m$ and
$\mathcal F$ have support $Y$ (I, 9.4.7). The restrictions
$v:\mathcal H_0\to\mathcal G^m$ and $w:\mathcal H_0\to\mathcal F$ of the
canonical projections of $\mathcal G^m\oplus\mathcal F$ are then
homomorphisms of coherent $\mathcal O_X$-modules which, in $W$ and in
$X-Y$, reduce to isomorphisms; in other words, the kernels and
cokernels of $v$ and $w$ have their supports in the closed set
$Y-(Y\cap W)$, which is distinct from $Y$. They therefore belong to $K'$; on the
other hand, we have $\mathcal G^m\in K'$ since $\mathcal G\in K'$ and
$K'$ is exact. We conclude successively, by the exactness of $K'$, that
$\mathcal H_0\in K'$ and then that $\mathcal F\in K'$. Q.E.D.
\end{enumerate}
\end{proof}

\begin{corollary}[3.1.3]
\label{III.3.1.3}
Suppose that the exact subset $K'$ of $K$ moreover has the property
that every coherent direct factor of a coherent $\mathcal O_X$-module
$\mathcal M\in K'$ also belongs to $K'$. Under these conditions, the
conclusion of (\hyperref[III.3.1.2]{3.1.2}) still holds when the
condition ``$\mathcal G_y$ is a $k(y)$-vector space of dimension $1$''
is replaced by $\mathcal G_y\neq0$ (which is equivalent to
$\operatorname{Supp}(\mathcal G)=Y$).
\end{corollary}

\begin{proof}
Indeed, the argument of (\hyperref[III.3.1.2]{3.1.2}) need only be
modified in case b); this time $\mathcal G_y$ is a $k(y)$-vector space of
dimension $q>0$, and we consequently have an
$\mathcal O_Y$-isomorphism $(\mathcal G_y)^m\xrightarrow{\sim}(\mathcal F_y)^q$; the end
of the argument of (\hyperref[III.3.1.2]{3.1.2}) then proves that
$\mathcal F^q\in K'$, and the additional hypothesis on $K'$ implies that
$\mathcal F\in K'$.
\end{proof}

\subsection{The finiteness theorem: the case of ordinary schemes}
\label{subsection:III.3.2}

\begin{theorem}[3.2.1]
\label{III.3.2.1}
Let $Y$ be a locally Noetherian prescheme and $f:X\to Y$ a proper
morphism. For every coherent $\mathcal O_X$-module $\mathcal F$, the
$\mathcal O_Y$-modules $R^qf_*(\mathcal F)$ are coherent for $q\geq0$.
\end{theorem}

\begin{proof}
The question being local on $Y$, we may suppose $Y$ Noetherian, hence $X$
Noetherian (I, 6.3.7). The coherent $\mathcal O_X$-modules $\mathcal F$ for
which the conclusion of Theorem (\hyperref[III.3.2.1]{3.2.1}) holds
form an exact subset $K'$ of the category $K$ of coherent
$\mathcal O_X$-modules.

\oldpage[III]{117}
Indeed, let
$0\to\mathcal F'\to\mathcal F\to\mathcal F''\to0$ be an exact sequence of
coherent $\mathcal O_X$-modules; suppose for example that $\mathcal F'$
and $\mathcal F''$ belong to $K'$; we have the exact cohomology sequence
\[
R^{q-1}f_*(\mathcal F'')\xrightarrow{\partial}R^qf_*(\mathcal F')
\longrightarrow R^qf_*(\mathcal F)\longrightarrow R^qf_*(\mathcal F'')
\xrightarrow{\partial}R^{q+1}f_*(\mathcal F')
\]
in which, by hypothesis, the four outer terms are coherent; the same is
therefore true of the middle term $R^qf_*(\mathcal F)$ by
(0\textsubscript{I}, 5.3.4 and 5.3.3). We show in the same way that
when $\mathcal F$ and $\mathcal F'$ (resp. $\mathcal F$ and
$\mathcal F''$) are in $K'$, the same is true of $\mathcal F''$ (resp.
$\mathcal F'$). Moreover, every coherent \emph{direct factor} $\mathcal F'$
of an $\mathcal O_X$-module $\mathcal F\in K'$ also belongs to $K'$:
indeed, $R^qf_*(\mathcal F')$ is then a direct factor of
$R^qf_*(\mathcal F)$ (G, II, 4.4.4), hence is of finite type, and since it is
quasi-coherent (\hyperref[III.1.4.10]{1.4.10}), it is coherent, $Y$
being Noetherian. By virtue of (\hyperref[III.3.1.3]{3.1.3}), we are
reduced to proving that when $X$ is \emph{irreducible} with generic point
$x$, there exists a coherent $\mathcal O_X$-module $\mathcal F$ belonging to
$K'$ such that $\mathcal F_x\neq0$: indeed, once this point is established,
we can apply it to every irreducible closed sub-prescheme $Y$ of $X$, since
if $j:Y\to X$ is the canonical injection, $f\circ j$ is proper (II, 5.4.2),
and if $\mathcal G$ is a coherent $\mathcal O_Y$-module with support $Y$,
$j_*(\mathcal G)$ is a coherent $\mathcal O_X$-module such that
\[
R^q(f\circ j)_*(\mathcal G)=R^qf_*(j_*(\mathcal G))
\]
(G, II, 4.9.1), so that we are indeed in the conditions for applying
(\hyperref[III.3.1.3]{3.1.3}).

Now, by virtue of Chow's lemma (II, 5.6.2), there exist an irreducible
prescheme $X'$ and a \emph{projective} and surjective morphism
$g:X'\to X$ such that $f\circ g:X'\to Y$ is \emph{projective}. There exists an
$\mathcal O_{X'}$-module $\mathcal L$ ample for $g$ (II, 5.3.1);
let us apply the fundamental theorem on projective morphisms
(\hyperref[III.2.2.1]{2.2.1}) to $g:X'\to X$ and to $\mathcal L$: there
then exists an integer $n$ such that
$\mathcal F=g_*(\mathcal O_{X'}(n))$ is a coherent $\mathcal O_X$-module
and $R^qg_*(\mathcal O_{X'}(n))=0$ for every $q>0$; moreover, since
\[
g^*(g_*(\mathcal O_{X'}(n)))\longrightarrow\mathcal O_{X'}(n)
\]
is surjective for $n$ sufficiently large (\hyperref[III.2.2.1]{2.2.1}), we
see that we may suppose that, at the generic point $x$ of $X$, we have
$\mathcal F_x\neq0$ (II, 3.4.7). On the other hand, since $f\circ g$ is
projective and $Y$ Noetherian, the
$R^p(f\circ g)_*(\mathcal O_{X'}(n))$ are \emph{coherent}
(\hyperref[III.2.2.1]{2.2.1}). This being so,
$R^*(f\circ g)_*(\mathcal O_{X'}(n))$ is the abutment of a Leray spectral
sequence whose $E_2$ term is given by
\[
E_2^{pq}=R^pf_*(R^qg_*(\mathcal O_{X'}(n)))~;
\]
what precedes shows that this spectral sequence is degenerate, and we then
know (0, 11.1.6) that $E_2^{p0}=R^pf_*(\mathcal F)$ is isomorphic to
$R^p(f\circ g)_*(\mathcal O_{X'}(n))$, which completes the proof.
\end{proof}

\begin{corollary}[3.2.2]
\label{III.3.2.2}
Let $Y$ be a locally Noetherian prescheme. For every proper morphism
$f:X\to Y$, the direct image under $f$ of every coherent $\mathcal O_X$-module
is a coherent $\mathcal O_Y$-module.
\end{corollary}

\begin{corollary}[3.2.3]
\label{III.3.2.3}
Let $A$ be a Noetherian ring and $X$ a scheme proper over $A$; for every
coherent $\mathcal O_X$-module $\mathcal F$, the $H^p(X,\mathcal F)$ are
$A$-modules of finite type, and there exists an integer $r>0$ such that for every
coherent $\mathcal O_X$-module $\mathcal F$ and every $p>r$,
$H^p(X,\mathcal F)=0$.
\end{corollary}

\begin{proof}
The second assertion has already been proved
(\hyperref[III.1.4.12]{1.4.12}); the first follows from the finiteness
theorem (\hyperref[III.3.2.1]{3.2.1}), taking into account
(\hyperref[III.1.4.11]{1.4.11}).
\end{proof}

In particular, if $X$ is a \emph{proper algebraic scheme} over a field
$k$, then, for every coherent $\mathcal O_X$-module $\mathcal F$, the
$H^p(X,\mathcal F)$ are $k$-vector spaces of \emph{finite
dimension}.

\begin{corollary}[3.2.4]
\label{III.3.2.4}
Let $Y$ be a locally Noetherian prescheme and $f:X\to Y$ a morphism of
finite type. For every coherent $\mathcal O_X$-module $\mathcal F$ whose
support is proper over $Y$ (II, 5.4.10), the $\mathcal O_Y$-modules
$R^qf_*(\mathcal F)$ are coherent.
\end{corollary}

\oldpage[III]{118}
\begin{proof}
The question being local on $Y$, we may suppose $Y$ Noetherian, and the
same then holds for $X$ (I, 6.3.7). By hypothesis, every closed sub-prescheme
$Z$ of $X$ whose underlying space is $\operatorname{Supp}(\mathcal F)$
is proper over $Y$; in other words, if $j:Z\to X$ is the canonical injection,
$f\circ j:Z\to Y$ is proper. Now, we may choose $Z$ such that
$\mathcal F=j_*(\mathcal G)$, where $\mathcal G=j^*(\mathcal F)$ is a
coherent $\mathcal O_Z$-module (I, 9.3.5); since we have
$R^qf_*(\mathcal F)=R^q(f\circ j)_*(\mathcal G)$ by
(\hyperref[III.1.3.4]{1.3.4}), the conclusion follows immediately from
(\hyperref[III.3.2.1]{3.2.1}).
\end{proof}

\subsection{Generalization of the finiteness theorem (ordinary schemes)}
\label{subsection:III.3.3}

\begin{proposition}[3.3.1]
\label{III.3.3.1}
Let $Y$ be a Noetherian prescheme, $\mathcal S$ an
$\mathcal O_Y$-algebra graded in positive degrees, quasi-coherent and of finite
type, $Y'=\operatorname{Proj}(\mathcal S)$, and $g:Y'\to Y$ the structural
morphism. Let $f:X\to Y$ be a proper morphism,
$\mathcal S'=f^*(\mathcal S)$, and
$\mathcal M=\bigoplus_{k\in\mathbb Z}\mathcal M_k$ a quasi-coherent graded
$\mathcal S'$-module of finite type. Then the
\[
R^pf_*(\mathcal M)=\bigoplus_{k\in\mathbb Z}R^pf_*(\mathcal M_k)
\]
are graded $\mathcal S$-modules of finite type for every $p$. Suppose
moreover that $\mathcal S$ is generated by $\mathcal S_1$; then, for
each $p\in\mathbb Z$, there exists an integer $k_p$ such that for every
$k\geq k_p$ and every $r\geq0$ we have
\[
R^pf_*(\mathcal M_{k+r})=\mathcal S_rR^pf_*(\mathcal M_k).
\tag{3.3.1.1}\label{III.3.3.1.1}
\]
\end{proposition}

\begin{proof}
The first assertion is identical to the statement of
(\hyperref[III.2.4.1]{2.4.1}, (i)), where we have simply replaced
``projective morphism'' by ``proper morphism''. Now, in the proof
of (\hyperref[III.2.4.1]{2.4.1}, (i)), the hypothesis on $f$ was used
only to show (with the notation of that proof) that
$R^pf'_*(\widetilde{\mathcal M})$ is a coherent $\mathcal O_{Y'}$-module.
Under the hypotheses of (\hyperref[III.3.3.1]{3.3.1}), $f'$ is proper
(II, 5.4.2, (iii)), so we can repeat without change the whole
proof of (\hyperref[III.2.4.1]{2.4.1}, (i)), thanks to the finiteness
theorem (\hyperref[III.3.2.1]{3.2.1}).

As for the second assertion, it suffices to remark that there is a
finite affine open covering $(U_i)$ of $Y$ such that the restrictions to
$U_i$ of the two sides of
(\hyperref[III.3.3.1.1]{3.3.1.1}) are equal for every
$k\geq k_{p,i}$ (II, 2.1.6, (ii)); it suffices to take for $k_p$ the
largest of the $k_{p,i}$.
\end{proof}

\begin{corollary}[3.3.2]
\label{III.3.3.2}
Let $A$ be a Noetherian ring, $\mathfrak m$ an ideal of $A$, $X$ a proper
$A$-scheme, and $\mathcal F$ a coherent $\mathcal O_X$-module. Then,
for every $p\geq0$, the direct sum
$\bigoplus_{k\geq0}H^p(X,\mathfrak m^k\mathcal F)$ is a module of finite type
over the ring $S=\bigoplus_{k\geq0}\mathfrak m^k$; in particular, there exists
an integer $k_p\geq0$ such that for every $k\geq k_p$ and every $r\geq0$ we have
\[
H^p(X,\mathfrak m^{k+r}\mathcal F)
=\mathfrak m^rH^p(X,\mathfrak m^k\mathcal F).
\tag{3.3.2.1}\label{III.3.3.2.1}
\]
\end{corollary}

\begin{proof}
It suffices to apply (\hyperref[III.3.3.1]{3.3.1}) with
$Y=\operatorname{Spec}(A)$, $\mathcal S=\widetilde S$,
$\mathcal M_k=\mathfrak m^k\mathcal F$, taking into account
(\hyperref[III.1.4.11]{1.4.11}).
\end{proof}

It should be recalled that the $S$-module structure of
$\bigoplus_{k\geq0}H^p(X,\mathfrak m^k\mathcal F)$ is obtained by considering,
for every $a\in\mathfrak m^r$, the map
\[
H^p(X,\mathfrak m^k\mathcal F)
\longrightarrow H^p(X,\mathfrak m^{k+r}\mathcal F)
\]
which arises, by passage to cohomology, from the multiplication
$\mathfrak m^k\mathcal F\to\mathfrak m^{k+r}\mathcal F$ defined by $a$
(\hyperref[III.2.4.1]{2.4.1}).

\oldpage[III]{119}
\subsection{The finiteness theorem: the case of formal schemes}
\label{subsection:III.3.4}

The results of this section (except for the definition
(\hyperref[III.3.4.1]{3.4.1})) will not be used in the rest of this
chapter.

\begin{env}[3.4.1]
\label{III.3.4.1}
Let $\mathfrak X$ and $\mathfrak S$ be two locally Noetherian formal
preschemes (I, 10.4.2) and $f:\mathfrak X\to\mathfrak S$ a morphism of
formal preschemes. We will say that $f$ is a \emph{proper} morphism if it
satisfies the following conditions:
\begin{enumerate}
\item[$1^\circ$] $f$ is a morphism of finite type (I, 10.13.3).
\item[$2^\circ$] If $\mathcal K$ is an ideal of definition of $\mathfrak S$
and if we set
$\mathcal J=f^*(\mathcal K)\mathcal O_{\mathfrak X}$,
$X_0=(\mathfrak X,\mathcal O_{\mathfrak X}/\mathcal J)$,
$S_0=(\mathfrak S,\mathcal O_{\mathfrak S}/\mathcal K)$, the morphism
$f_0:X_0\to S_0$ deduced from $f$ (I, 10.5.6) is proper.
\end{enumerate}

It is immediate that this definition does not depend on the ideal of
definition $\mathcal K$ of $\mathfrak S$ under consideration; indeed, if
$\mathcal K'$ is a second ideal of definition such that
$\mathcal K'\subset\mathcal K$, and if we set
$\mathcal J'=f^*(\mathcal K')\mathcal O_{\mathfrak X}$,
$X'_0=(\mathfrak X,\mathcal O_{\mathfrak X}/\mathcal J')$,
$S'_0=(\mathfrak S,\mathcal O_{\mathfrak S}/\mathcal K')$, the morphism
$f'_0:X'_0\to S'_0$ deduced from $f$ is such that the diagram
\[
\xymatrix{
X_0\ar[r]^{f_0}\ar[d]_i & S_0\ar[d]^j\\
X'_0\ar[r]_{f'_0} & S'_0
}
\]
is commutative, $i$ and $j$ being surjective immersions; it therefore amounts
to the same to say that $f_0$ or that $f'_0$ is proper, by virtue of (II, 5.4.5).

Note that, for every $n\geq0$, if we set
$X_n=(\mathfrak X,\mathcal O_{\mathfrak X}/\mathcal J^{n+1})$,
$S_n=(\mathfrak S,\mathcal O_{\mathfrak S}/\mathcal K^{n+1})$, the morphism
$f_n:X_n\to S_n$ deduced from $f$ (I, 10.5.6) is proper for every $n$ as soon
as it is so for $n=0$ (II, 5.4.6).

If $g:Y\to Z$ is a proper morphism of locally Noetherian ordinary
preschemes, $Z'$ a closed subset of $Z$, and $Y'$ a closed subset of $Y$
such that $g(Y')\subset Z'$, the extension
$\widehat g:Y_{/Y'}\to Z_{/Z'}$ of $g$ to the completions (I, 10.9.1) is a
proper morphism of formal preschemes, as follows from the definition and
from (II, 5.4.5).

Let $\mathfrak X$ and $\mathfrak S$ be two locally Noetherian formal
preschemes and $f:\mathfrak X\to\mathfrak S$ a morphism \emph{of finite type}
(I, 10.13.3); the notation being the same as above, we say that a
subset $Z$ of the underlying space of $\mathfrak X$ is \emph{proper over}
$\mathfrak S$ (or proper for $f$) if, considered as a subset of $X_0$, $Z$
is \emph{proper over $S_0$} (II, 5.4.10). All the properties of proper
subsets of ordinary preschemes stated in (II, 5.4.10) remain valid
for proper subsets of formal preschemes, as follows immediately from the
definitions.
\end{env}

\begin{theorem}[3.4.2]
\label{III.3.4.2}
Let $\mathfrak X$, $\mathfrak Y$ be two locally Noetherian formal
preschemes and $f:\mathfrak X\to\mathfrak Y$ a proper morphism. For every
coherent $\mathcal O_{\mathfrak X}$-module $\mathcal F$, the
$\mathcal O_{\mathfrak Y}$-modules $R^qf_*(\mathcal F)$ are coherent for
every $q\geq0$.
\end{theorem}

Let $\mathcal J$ be an ideal of definition of $\mathfrak Y$ and
$\mathcal K=f^*(\mathcal J)\mathcal O_{\mathfrak X}$, and consider the
$\mathcal O_{\mathfrak X}$-modules
\[
\mathcal F_k=\mathcal F\otimes_{\mathcal O_{\mathfrak Y}}
(\mathcal O_{\mathfrak Y}/\mathcal J^{k+1})
=\mathcal F/\mathcal K^{k+1}\mathcal F\qquad(k\geq0)
\tag{3.4.2.1}\label{III.3.4.2.1}
\]
which evidently form a \emph{projective system} of topological
$\mathcal O_{\mathfrak X}$-modules such that

\oldpage[III]{120}
\[
\mathcal F=\varprojlim_k\mathcal F_k
\]
by (I, 10.11.3). On the other hand, it will follow from
(\hyperref[III.3.4.2]{3.4.2}) that each of the $R^qf_*(\mathcal F)$, being
coherent, is naturally endowed with a structure of topological
$\mathcal O_{\mathfrak Y}$-module (I, 10.11.6), and the same holds
for the $R^qf_*(\mathcal F_k)$. To the canonical homomorphisms
$\mathcal F\to\mathcal F_k=\mathcal F/\mathcal K^{k+1}\mathcal F$
there canonically correspond homomorphisms
\[
R^qf_*(\mathcal F)\longrightarrow R^qf_*(\mathcal F_k)
\]
which are necessarily continuous for the preceding structures of topological
$\mathcal O_{\mathfrak Y}$-module (I, 10.11.6), and
form a projective system, giving in the limit a canonical functorial
homomorphism
\[
R^qf_*(\mathcal F)\longrightarrow\varprojlim_kR^qf_*(\mathcal F_k)
\tag{3.4.2.2}\label{III.3.4.2.2}
\]
which will be a continuous homomorphism of topological
$\mathcal O_{\mathfrak Y}$-modules. We will prove, at the same time as
(\hyperref[III.3.4.2]{3.4.2}), the

\begin{corollary}[3.4.3]
\label{III.3.4.3}
Each of the homomorphisms (\hyperref[III.3.4.2.2]{3.4.2.2}) is a
topological isomorphism. Moreover, if $\mathfrak Y$ is Noetherian, the
projective system
$(R^qf_*(\mathcal F/\mathcal K^{k+1}\mathcal F))_{k\geq0}$ satisfies
condition (ML) (0, 13.1.1).
\end{corollary}

We will begin by establishing (\hyperref[III.3.4.2]{3.4.2}) and
(\hyperref[III.3.4.3]{3.4.3}) when $\mathfrak Y$ is a Noetherian affine
formal scheme (I, 10.4.1):

\begin{corollary}[3.4.4]
\label{III.3.4.4}
Under the hypotheses of (\hyperref[III.3.4.2]{3.4.2}), suppose moreover
that $\mathfrak Y=\operatorname{Spf}(A)$, where $A$ is a Noetherian adic
ring. Let $\mathfrak J$ be an ideal of definition of $A$, and set
$\mathcal F_k=\mathcal F/\mathfrak J^{k+1}\mathcal F$ for $k\geq0$. Then
the $H^n(\mathfrak X,\mathcal F)$ are $A$-modules of finite type; the
projective system $(H^n(\mathfrak X,\mathcal F_k))_{k\geq0}$ satisfies
condition (ML) for every $n$; if we set
\[
N_{n,k}=\operatorname{Ker}\bigl(H^n(\mathfrak X,\mathcal F)
\longrightarrow H^n(\mathfrak X,\mathcal F_k)\bigr)
\tag{3.4.4.1}\label{III.3.4.4.1}
\]
(also equal to
$\operatorname{Im}(H^n(\mathfrak X,\mathfrak J^{k+1}\mathcal F)
\to H^n(\mathfrak X,\mathcal F))$ by the exact cohomology sequence), the
$N_{n,k}$ define on $H^n(\mathfrak X,\mathcal F)$ a
$\mathfrak J$-good filtration (0, 13.7.7); finally, the canonical homomorphism
\[
H^n(\mathfrak X,\mathcal F)
\longrightarrow\varprojlim_kH^n(\mathfrak X,\mathcal F_k)
\tag{3.4.4.2}\label{III.3.4.4.2}
\]
is a topological isomorphism for every $n$ (the left-hand side being endowed
with the $\mathfrak J$-adic topology, and the $H^n(\mathfrak X,\mathcal F_k)$ with
the discrete topology).
\end{corollary}

Set
\[
S=\operatorname{gr}(A)=\bigoplus_{k\geq0}
\mathfrak J^k/\mathfrak J^{k+1},\qquad
\mathcal M=\operatorname{gr}(\mathcal F)=\bigoplus_{k\geq0}
\mathfrak J^k\mathcal F/\mathfrak J^{k+1}\mathcal F.
\tag{3.4.4.3}\label{III.3.4.4.3}
\]
We know that $\mathfrak J^\Delta$ is an ideal of definition of $\mathfrak Y$
(I, 10.3.1); let
$\mathcal K=f^*(\mathfrak J^\Delta)\mathcal O_{\mathfrak X}$,
$X_0=(\mathfrak X,\mathcal O_{\mathfrak X}/\mathcal K)$,
$Y_0=(\mathfrak Y,\mathcal O_{\mathfrak Y}/\mathfrak J^\Delta)
=\operatorname{Spec}(A_0)$ with $A_0=A/\mathfrak J$. It is clear that the
$\mathcal M_k=\mathfrak J^k\mathcal F/\mathfrak J^{k+1}\mathcal F$ are
coherent $\mathcal O_{X_0}$-modules (I, 10.11.3). Consider on the other hand the
quasi-coherent graded $\mathcal O_{X_0}$-algebra
\[
\mathcal S=\mathcal O_{X_0}\otimes_{A_0}S.
\tag{3.4.4.4}\label{III.3.4.4.4}
\]
Multiplication defines a surjective homomorphism of graded algebras
\[
\mathcal S\longrightarrow
\operatorname{gr}(\mathcal O_{\mathfrak X})
=\bigoplus_{k\geq0}\mathcal K^k/\mathcal K^{k+1},
\]
by means of which $\mathcal M$ is a graded $\mathcal S$-module.

The hypothesis that $\mathcal F$ is an $\mathcal O_{\mathfrak X}$-module of
finite type implies first of all that $\mathcal M$ is

\oldpage[III]{121}
a graded $\mathcal S$-module \emph{of finite type}. Indeed, the question is
local on $\mathfrak X$, and to treat it we may therefore suppose that
$\mathfrak X=\operatorname{Spf}(B)$, where $B$ is a Noetherian adic ring,
and $\mathcal F=N^\Delta$, where $N$ is a $B$-module of finite type
(I, 10.10.5). Set $B_0=B/\mathfrak J B$ and $N_0=N/\mathfrak JN$.
Then $X_0=\operatorname{Spec}(B_0)$, and the
quasi-coherent $\mathcal O_{X_0}$-modules $\mathcal S$ and $\mathcal M$
are associated respectively to
\[
S'=\bigoplus_{k\geq0}
((\mathfrak J^k/\mathfrak J^{k+1})\otimes_{A_0}B_0),\qquad
M'=\bigoplus_{k\geq0}\mathfrak J^kN/\mathfrak J^{k+1}N.
\]
For every $k\geq0$, multiplication induces a surjective map
\[
(\mathfrak J^k/\mathfrak J^{k+1})\otimes_{A_0}N_0
\longrightarrow\mathfrak J^kN/\mathfrak J^{k+1}N,
\qquad \overline a\otimes\overline n\longmapsto\overline{an}.
\]
These maps define a surjective homomorphism of graded
$S'$-modules $S'\otimes_{B_0}N_0\longrightarrow M'$.
Since $N_0$ is a $B_0$-module of finite type, $M'$ is therefore an
$S'$-module of finite type, whence our assertion (I, 1.3.13).

Since the morphism $f_0:X_0\to Y_0$ is \emph{proper} by hypothesis, we can
apply (\hyperref[III.3.3.1]{3.3.1}) to $\mathcal S$, $\mathcal M$ and the
morphism $f_0$; taking into account (\hyperref[III.1.4.11]{1.4.11}), we
conclude that for \emph{every $n\geq0$},
$\bigoplus_{k\geq0}H^n(X_0,\mathcal M_k)$ is a graded $S$-module \emph{of
finite type}. This proves that condition $(\mathrm F_n)$ of (0, 13.7.7) is
satisfied for \emph{every $n\geq0$} when we consider the strict projective
system $(\mathcal F/\mathfrak J^k\mathcal F)_{k\geq0}$ of sheaves of
abelian groups on $X_0$, each endowed with its natural structure of
``filtered $A$-module''. We can therefore apply (0, 13.7.7), which proves that:
\begin{enumerate}
\item[$1^\circ$] The projective system
$(H^n(\mathfrak X,\mathcal F_k))_{k\geq0}$ satisfies condition (ML).
\item[$2^\circ$] If
$H'^{,n}=\varprojlim_kH^n(\mathfrak X,\mathcal F_k)$, then $H'^{,n}$ is an
$A$-module of finite type.
\item[$3^\circ$] The filtration defined on $H'^{,n}$ by the kernels of the
canonical homomorphisms $H'^{,n}\to H^n(\mathfrak X,\mathcal F_k)$ is
$\mathfrak J$-good.
\end{enumerate}

Note on the other hand that if we set
$X_k=(\mathfrak X,\mathcal O_{\mathfrak X}/\mathcal K^{k+1})$, then
$\mathcal F_k$ is a coherent $\mathcal O_{X_k}$-module (I, 10.11.3), and if
$U$ is an affine open subset of $X_0$, then $U$ is also an affine open subset of
each of the $X_k$ (I, 5.1.9), hence $H^n(U,\mathcal F_k)=0$ for every $n>0$ and
every $k$ (\hyperref[III.1.3.1]{1.3.1}), and
$H^0(U,\mathcal F_k)\to H^0(U,\mathcal F_h)$ is surjective for $h\leq k$
(I, 1.3.9). We are therefore in the conditions of (0, 13.3.2), and the application
of (0, 13.3.1) proves that $H'^{,n}$ is canonically identified with
$H^n(\mathfrak X,\varprojlim_k\mathcal F_k)=H^n(\mathfrak X,\mathcal F)$;
this completes the proof of (\hyperref[III.3.4.4]{3.4.4}).

\begin{env}[3.4.5]
\label{III.3.4.5}
Let us now return to the proof of (\hyperref[III.3.4.2]{3.4.2})
and (\hyperref[III.3.4.3]{3.4.3}). Let us first prove these propositions in
the case $\mathfrak Y=\operatorname{Spf}(A)$ considered in
(\hyperref[III.3.4.4]{3.4.4}); to do this, for every $g\in A$, let us apply
(\hyperref[III.3.4.4]{3.4.4}) to the Noetherian affine formal scheme induced
on the open subset $\mathfrak Y_g=\mathfrak D(g)$ of $\mathfrak Y$, which is equal to
$\operatorname{Spf}(A_{\{g\}})$, and to the formal prescheme induced by
$\mathfrak X$ on $f^{-1}(\mathfrak Y_g)$; note that $\mathfrak Y_g$ is
also an affine open subset of the prescheme
$Y_k=(\mathfrak Y,\mathcal O_{\mathfrak Y}/(\mathfrak J^\Delta)^{k+1})$,
and since $\mathcal F_k$ is a coherent $\mathcal O_{X_k}$-module, we have
\[
H^n(f^{-1}(\mathfrak Y_g),\mathcal F_k)
=\Gamma(\mathfrak Y_g,R^nf_*(\mathcal F_k))
\]
for every $k\geq0$ by virtue of
(\hyperref[III.1.4.11]{1.4.11}). The canonical homomorphism
\[
H^n(f^{-1}(\mathfrak Y_g),\mathcal F)
\longrightarrow\varprojlim_k\Gamma(\mathfrak Y_g,R^nf_*(\mathcal F_k))
\]
is an isomorphism; but we have (0\textsubscript{I}, 3.2.6)
\[
\varprojlim_k\Gamma(\mathfrak Y_g,R^nf_*(\mathcal F_k))
=\Gamma(\mathfrak Y_g,\varprojlim_kR^nf_*(\mathcal F_k))
\]

\oldpage[III]{122}
and since the sheaf $R^nf_*(\mathcal F)$ is associated to the presheaf
$\mathfrak Y_g\mapsto H^n(f^{-1}(\mathfrak Y_g),\mathcal F)$ on the
$\mathfrak Y_g$ (0\textsubscript{I}, 3.2.1), we have indeed shown that
the homomorphism (\hyperref[III.3.4.2.2]{3.4.2.2}) is bijective. Let us
next prove that $R^nf_*(\mathcal F)$ is a coherent $\mathcal O_{\mathfrak Y}$-module,
and more precisely that we have
\[
R^nf_*(\mathcal F)=(H^n(\mathfrak X,\mathcal F))^\Delta.
\tag{3.4.5.1}\label{III.3.4.5.1}
\]
With the preceding notation, we have, since $\mathcal F_k$ is a coherent
$\mathcal O_{X_k}$-module (\hyperref[III.1.4.13]{1.4.13}),
\[
\Gamma(\mathfrak Y_g,R^nf_*(\mathcal F_k))
=(\Gamma(\mathfrak Y,R^nf_*(\mathcal F_k)))_g
=(H^n(\mathfrak X,\mathcal F_k))_g.
\]
Now, the $H^n(\mathfrak X,\mathcal F_k)$ form a projective system
satisfying (ML), and their projective limit $H^n(\mathfrak X,\mathcal F)$ is
an $A$-module of finite type. We conclude from this (0, 13.7.8) that we have
\[
\varprojlim_k(H^n(\mathfrak X,\mathcal F_k))_g
=H^n(\mathfrak X,\mathcal F)\otimes_AA_{\{g\}}
=\Gamma(\mathfrak Y_g,(H^n(\mathfrak X,\mathcal F))^\Delta)
\]
taking into account (I, 10.10.8) applied to $A$ and $A_{\{g\}}$; this proves
(\hyperref[III.3.4.5.1]{3.4.5.1}), since
\[
\Gamma(\mathfrak Y_g,R^nf_*(\mathcal F))
=\varprojlim_k\Gamma(\mathfrak Y_g,R^nf_*(\mathcal F_k)).
\]
Since (\hyperref[III.3.4.2.2]{3.4.2.2}) is then an isomorphism of coherent
$\mathcal O_{\mathfrak Y}$-modules, it is necessarily a
\emph{topological} isomorphism (I, 10.11.6). Finally, it follows from the relations
$R^nf_*(\mathcal F_k)=(H^n(\mathfrak X,\mathcal F_k))^\Delta$ that the projective
system $(R^nf_*(\mathcal F_k))_{k\geq0}$ satisfies (ML) (I, 10.10.2).

Once (\hyperref[III.3.4.2]{3.4.2}) and
(\hyperref[III.3.4.3]{3.4.3}) have been proved in the case where the formal
prescheme $\mathfrak Y$ is Noetherian affine, it is immediate to pass from there to the
general case for (\hyperref[III.3.4.2]{3.4.2}) and the first assertion of
(\hyperref[III.3.4.3]{3.4.3}), which are local on $\mathfrak Y$. As for
the second assertion of (\hyperref[III.3.4.3]{3.4.3}), it suffices,
$\mathfrak Y$ being Noetherian, to cover it by finitely many Noetherian
affine open subsets $U_i$ and to remark that the restrictions of the projective
system $(R^qf_*(\mathcal F_k))$ to each of the $U_i$ satisfy (ML).
\end{env}

We have moreover proved along the way:

\begin{corollary}[3.4.6]
\label{III.3.4.6}
Under the hypotheses of (\hyperref[III.3.4.4]{3.4.4}), the canonical
homomorphism
\[
H^q(\mathfrak X,\mathcal F)
\longrightarrow\Gamma(\mathfrak Y,R^qf_*(\mathcal F))
\tag{3.4.6.1}\label{III.3.4.6.1}
\]
is bijective.
\end{corollary}

